% Book 4, Year 2, chapter 16 -- Sockets.
% Provenance (never printed): sources/book4/BRIEFS.md, chapter 16.
\chapter{Sockets}\label{ch:b2:sockets}

Two programs on one laptop, no cable between them, and still a conversation:

\omoutput{code/bachelor-2/16-sockets/outputs/echo_pair.txt}{A server and a client, two \omterm{def:b2:processes-and-system-calls:process}{processes}, one
connection.}\label{out:b2:sockets:hook}

The client sent 13 bytes; the server received them, saw that they came from the address \code{127.0.0.1}, and sent
them back. The bytes went down through the operating system's network code and came back up, exactly as they would
between two machines on two continents, except that they never left this one. The tool is the \emph{\omterm{def:b2:sockets:socket}{socket}}, a \omterm{def:b2:files-pipes-signals:descriptor}{file
descriptor} for a conversation over the network. This chapter writes clients and servers with \omterm{def:b2:sockets:socket}{sockets}, shows the two
traps every network program meets (a stream has no message boundaries; a number has a byte order), serves many
clients at once in three ways, and measures them.

\begin{recall}
Book 2: a web browser is a client that sends requests to a server, which answers; a machine on a network has an
address, and a \emph{\omterm{def:b2:sockets:port}{port number}} names one program on it, so that one machine can run several servers. Book 3, ch.~6:
the byte order of a machine says in which order the bytes of a number lie in memory. \cref{ch:b2:files-pipes-signals}:
descriptors, \code{read} and \code{write}, and end of file on a \omterm{def:b2:files-pipes-signals:pipe}{pipe}. \cref{ch:b2:threads-and-synchronisation}:
\omterm{def:b2:threads-and-synchronisation:thread}{threads}, \omterm{def:b2:threads-and-synchronisation:lock}{locks}, and the cost of a \omterm{def:b2:threads-and-synchronisation:thread}{thread}.
\end{recall}

\section{The socket interface}\label{sec:b2:sockets:interface}

\begin{definition}[Socket, socket address]\label{def:b2:sockets:socket}
A \emph{socket}\index{socket} is an endpoint of communication, created by \code{socket(domain, type, 0)}, which returns
a \omterm{def:b2:files-pipes-signals:descriptor}{file descriptor}. A \emph{socket address}\index{socket address} names an endpoint: for the Internet domain
(\code{AF_INET}), an IP address and a \omterm{def:b2:sockets:port}{port number}, held in a \code{struct sockaddr_in}.
\end{definition}

\begin{definition}[Port number]\label{def:b2:sockets:port}
A \emph{port number}\index{port number} is a 16-bit number that tells the endpoints of one machine apart: the system
delivers each arriving segment or datagram to the \omterm{def:b2:sockets:socket}{socket} bound to its destination address and port.
\end{definition}

\begin{definition}[Stream socket, datagram socket]\label{def:b2:sockets:kinds}
A \emph{stream socket}\index{stream socket} (\code{SOCK_STREAM}, the TCP protocol) carries a reliable, ordered sequence
of bytes between two connected endpoints; a \emph{datagram socket}\index{datagram socket} (\code{SOCK_DGRAM}, the UDP
protocol) carries separate messages, each addressed on its own.
\end{definition}

% ledger: b2:loopback-block
\begin{definition}[Loopback address]\label{def:b2:sockets:loopback}
A \emph{loopback address}\index{loopback address} is an address of the block \code{127.0.0.0/8}, usually
\code{127.0.0.1}: a packet sent to it never leaves the machine.
\end{definition}

Every program of this chapter uses it, so nothing it does is visible from outside. The
transport chapter (\cref{ch:b2:transport-layer}) explains what TCP and UDP do underneath; here we use them.

A server and a client make different calls, in a fixed order:

\begin{method}[A TCP server and client]\label{met:b2:sockets:server-skeleton}
Server: \code{socket}; \code{bind} it to its address and port; \code{listen}, which turns it into a \omterm{def:b2:sockets:listening}{listening
socket}; then repeatedly \code{accept}, which waits for a client and returns a \emph{new} \omterm{def:b2:sockets:socket}{socket} for that connection;
\code{recv} and \code{send} on that \omterm{def:b2:sockets:socket}{socket}; \code{close} it. Client: \code{socket}; \code{connect} to the server's
address and port (the system picks the client's own port); \code{send} and \code{recv}; \code{close}. Every call
returns $-1$ on failure, which a real program checks.
\end{method}

\begin{definition}[Listening socket, connected socket]\label{def:b2:sockets:listening}
A \emph{listening socket}\index{listening socket} is a \omterm{def:b2:sockets:kinds}{stream socket} on which \code{listen} was called: it carries no
data, only connection requests, and \code{accept} turns each request into a \emph{connected socket}\index{connected
socket}, one per client, which carries that connection's bytes.
\end{definition}

\omcode{code/bachelor-2/16-sockets/listings/echo_pair.c}{42}{65}{The server's \omterm{def:b2:sockets:socket}{socket}: an address, then
\code{listen}; the client runs in a \omterm{def:b2:processes-and-system-calls:tree}{child process}}

\omcode{code/bachelor-2/16-sockets/listings/echo_pair.c}{14}{40}{The client connects and sends; the server accepts,
receives and echoes}\label{lst:b2:sockets:echo}

\begin{omfigure}
\begin{tikzpicture}[font=\scriptsize, c/.style={draw, rounded corners=2pt, fill=omDef!10, minimum width=1.7cm,
  minimum height=5mm, inner sep=2pt}, s/.style={draw, rounded corners=2pt, fill=omThm!12, minimum width=1.7cm,
  minimum height=5mm, inner sep=2pt}, w/.style={-{Stealth[length=1.8mm]}}]
\node[font=\small\bfseries] at (0,0.6) {server}; \node[font=\small\bfseries] at (7.0,0.6) {client};
\node[s] (s1) at (0,0) {\code{socket}}; \node[s] (s2) at (0,-0.7) {\code{bind}};
\node[s] (s3) at (0,-1.4) {\code{listen}}; \node[s] (s4) at (0,-2.1) {\code{accept}};
\node[s] (s5) at (0,-3.5) {\code{recv}}; \node[s] (s6) at (0,-4.2) {\code{send}}; \node[s] (s7) at (0,-4.9)
{\code{close}};
\node[c] (c1) at (7.0,-1.4) {\code{socket}}; \node[c] (c2) at (7.0,-2.1) {\code{connect}};
\node[c] (c3) at (7.0,-2.8) {\code{send}}; \node[c] (c4) at (7.0,-4.2) {\code{recv}}; \node[c] (c5) at (7.0,-4.9)
{\code{close}};
\foreach \a/\b in {s1/s2,s2/s3,s3/s4,s4/s5,s5/s6,s6/s7,c1/c2,c2/c3,c3/c4,c4/c5} \draw[w] (\a) -- (\b);
\draw[w, omThm, thick] (c2.west) -- node[above, sloped] {connection set up} (s4.east);
\draw[w, omThm, thick] (c3.west) -- node[above, sloped] {13 bytes} (s5.east);
\draw[w, omThm, thick] (s6.east) -- node[above] {13 bytes back} (c4.west);
\node[anchor=east, align=right] at (-1.0,-2.1) {\code{accept} waits for a client,\\ then returns a\\ new socket};
\end{tikzpicture}
\par\smallskip{\small The calls of \cref{lst:b2:sockets:echo} in time, top to bottom. \code{accept} returns when a
client's \code{connect} arrives; the data then flows on the new \omterm{def:b2:sockets:socket}{socket}, while the listening one can wait for the
next client.}
\end{omfigure}

The client is forked after \code{listen}, so its \code{connect} cannot arrive before the server is ready. The server
prints and \omterm{def:b2:files-pipes-signals:buffering}{flushes} before it answers, and the client prints only after the answer arrives, which fixes the order of
the two lines in \cref{out:b2:sockets:hook}.

\section{Streams have no message boundaries}\label{sec:b2:sockets:streams}

The client sent 13 bytes in one \code{send}, and the server received 13 bytes in one \code{recv}. That was luck. A
\omterm{def:b2:sockets:kinds}{stream socket} carries bytes, not messages: the receiver gets the bytes in order, but a \code{recv} may return any
number of them, from one to all that have arrived.

\omoutput{code/bachelor-2/16-sockets/outputs/partial_recv.txt}{One \code{send} of 1 MiB, received by several
\code{recv} calls.}

\omoutput{code/bachelor-2/16-sockets/outputs/boundaries.txt}{Three sends over each kind of \omterm{def:b2:sockets:socket}{socket}: the stream merges
them, datagrams keep them apart.}\label{out:b2:sockets:boundaries}

Over TCP, the three sends ``ab'', ``cd'' and ``efgh'' arrive as eight bytes that one \code{recv} returns together; it
could as well have returned them in five pieces. Over UDP, each \code{sendto} is one datagram, and each
\code{recvfrom} returns exactly one. A datagram can be lost, though, and one longer than the buffer is cut
(\cref{exo:b2:sockets:5}); TCP loses nothing and keeps the order, but leaves the boundaries to the program.

\begin{omfigure}
\begin{tikzpicture}[font=\scriptsize, b/.style={draw, minimum height=5mm, inner sep=2pt, fill=omDef!12},
  w/.style={-{Stealth[length=1.8mm]}}]
\node[anchor=east] at (-0.2,0.6) {sent}; \node[anchor=east] at (-0.2,-0.3) {TCP receives};
\node[anchor=east] at (-0.2,-1.2) {UDP receives};
\node[b] at (0.4,0.6) {\code{ab}}; \node[b] at (1.4,0.6) {\code{cd}}; \node[b] at (2.6,0.6) {\code{efgh}};
\node[b, fill=omThm!12, minimum width=2.6cm] at (1.5,-0.3) {\code{abcdefgh}};
\node[font=\scriptsize, anchor=west] at (3.1,-0.3) {one \code{recv}; or \code{abc} + \code{defgh}; or eight of one
byte \dots};
\node[b, fill=omThm!12] at (0.4,-1.2) {\code{ab}}; \node[b, fill=omThm!12] at (1.4,-1.2) {\code{cd}};
\node[b, fill=omThm!12] at (2.6,-1.2) {\code{efgh}};
\node[font=\scriptsize, anchor=west] at (3.1,-1.2) {three \code{recvfrom}, one datagram each (or a datagram lost)};
\end{tikzpicture}
\par\smallskip{\small A stream against datagrams. A stream keeps the bytes and their order, not the cuts between
sends; datagrams keep the cuts, not the guarantee.}
\end{omfigure}

So a program that needs $n$ bytes must loop:

\omcode{code/bachelor-2/16-sockets/listings/net.c}{56}{68}{Receiving exactly $n$ bytes, or what is left before the
end of the stream}\label{lst:b2:sockets:recv-all}

\begin{proposition}[Receiving everything]\label{prop:b2:sockets:recv-all}
If no call fails, \code{recv_all(fd, buf, n)} returns $n$ with the next $n$ bytes of the stream in \code{buf}, or,
if the stream ends first, returns the number $k < n$ of bytes that remained, with those $k$ bytes in \code{buf}.
\end{proposition}

\begin{proof}
The loop keeps the invariant: \code{buf[0..got)} holds the next \code{got} bytes of the stream, in order, and
$\text{got} \le n$. It holds at the start ($\text{got} = 0$). A \code{recv} that returns $k > 0$ copies the next $k$
bytes, at most $n - \text{got}$ of them, to \code{buf + got}, and \code{got} grows by $k$: the invariant still holds.
Each turn either increases \code{got} or ends the loop, so the loop ends: either $\text{got} = n$, or \code{recv}
returned 0, which on a stream means that the other side has closed and no byte remains.
\end{proof}

\code{send} has the same weakness in the other direction: it may accept fewer bytes than given, and
\code{send_all} loops in the same way.

\section{Byte order and framing}\label{sec:b2:sockets:framing}

% ledger: b2:network-byte-order
Numbers on the network have a byte order too.

\begin{definition}[Network byte order]\label{def:b2:sockets:network-order}
\emph{Network byte order}\index{network byte order} is big-endian: the Internet protocols send every number of a
header with its most significant byte first.
\end{definition}

This laptop is little-endian, so every number that goes into a protocol
header or a message must be converted: \code{htons} and \code{htonl} (host to network, short and long) and
\code{ntohs}, \code{ntohl} back.

\omoutput{code/bachelor-2/16-sockets/outputs/byteorder.txt}{The same numbers in memory and in network
order.}\label{out:b2:sockets:bytes}

A \omterm{def:b2:sockets:port}{port number} goes into \code{sin_port} in network order. Forgetting \code{htons} does not fail: the server listens
on the port whose bytes are those of 44164 reversed, and the client, which asks for 44164, finds nobody:

\omoutput{code/bachelor-2/16-sockets/outputs/wrong_port.txt}{A forgotten \code{htons}
(\cref{exo:b2:sockets:2}).}

Since the stream does not keep message boundaries, a protocol must mark them itself. There are two ways. A
\emph{delimiter} ends each message, as the newline ends a line in many text protocols; the message must then never
contain the delimiter, or must escape it. A \emph{length prefix} puts the length first; the message may then contain
anything.

\begin{definition}[Length-prefixed framing]\label{def:b2:sockets:framing}
In \emph{length-prefixed framing}\index{length-prefixed framing}, a message of $n$ bytes is sent as 4 bytes holding
$n$ in \omterm{def:b2:sockets:network-order}{network byte order}, followed by the $n$ bytes. Decoding a buffer reads 4 bytes as a length $n$, then takes the
next $n$ bytes as the message, and starts again after it.
\end{definition}

\omcode{code/bachelor-2/16-sockets/listings/net.c}{70}{89}{Framing and unframing: a whole message, or nothing
yet}\label{lst:b2:sockets:frame}

\begin{omfigure}
\begin{bytefield}[bitwidth=2.6em]{9}
\bitheader{0-8} \\
\bitbox{1}{\scriptsize 00} & \bitbox{1}{\scriptsize 00} & \bitbox{1}{\scriptsize 00} & \bitbox{1}{\scriptsize 05}
& \bitbox{1}{\scriptsize\code{w}} & \bitbox{1}{\scriptsize\code{h}} & \bitbox{1}{\scriptsize\code{o}}
& \bitbox{1}{\scriptsize\textvisiblespace} & \bitbox{1}{\scriptsize\code{1}}
\end{bytefield}
\par\smallskip{\small The request \code{who 1} framed, byte by byte (offsets above): the length 5 in four big-endian
bytes, then the five bytes of the message. The receiver reads the four bytes, then exactly five more; whatever follows is the next message.}
\end{omfigure}

\begin{proposition}[Unique decoding]\label{prop:b2:sockets:prefix-free}
Let $w$ be the concatenation of the framings of messages $m_1, \dots, m_k$. Decoding $w$ from its first byte yields
exactly $m_1, \dots, m_k$, in that order, and uses all of $w$.
\end{proposition}

The proof, by induction on $k$, is \cref{exo:b2:sockets:8}. A delimiter gives the same guarantee only if no message
contains it.

\section{A directory service}\label{sec:b2:sockets:directory}

The town keeps its people and streets in two tables. A small service answers two requests over TCP, each framed as
above: \code{who 7} gives a person's name, year of birth and street; \code{street Mill Lane} counts the people who
live there. Every connection carries one request and one answer. The client is a short Python program, whose
\code{struct.pack("!I", n)} writes $n$ in four big-endian bytes (the \code{!} is network order):

\omcode{code/bachelor-2/16-sockets/listings/directory_client.py}{7}{22}{A client in Python: the same framing, the same
loop}

\omoutput{code/bachelor-2/16-sockets/outputs/directory_demo.txt}{The server in C, its client in
Python.}\label{out:b2:sockets:directory}

The server's \code{&} runs it in the background; it serves four connections and ends, and \code{wait} waits for it. A
C server and a Python client agree because the protocol fixes every byte: the length's size and order, the text's
encoding. Neither side's language matters.

\section{Many clients at once}\label{sec:b2:sockets:concurrency}

A server that handles one client from \code{accept} to \code{close} makes every other client wait, even when the
first is slow to send. The calls it makes wait.

\begin{definition}[Blocking call, non-blocking descriptor]\label{def:b2:sockets:blocking}
A \emph{blocking call}\index{blocking call} waits until it can complete: \code{accept} until a client connects,
\code{recv} until a byte has arrived. On a descriptor made \emph{non-blocking}\index{non-blocking descriptor} (flag
\code{O_NONBLOCK}), the same calls return at once with the error \code{EAGAIN} instead of waiting.
\end{definition}

There are three classic ways to serve several clients at once.

\omcode{code/bachelor-2/16-sockets/listings/servers.c}{27}{42}{A \omterm{def:b2:processes-and-system-calls:process}{process} per connection: the child serves, the
parent accepts again}

The child inherits the \omterm{def:b2:sockets:listening}{connected socket} and the parent closes its copy (the connection ends only when both have, as
for a \omterm{def:b2:files-pipes-signals:pipe}{pipe}). The children are reaped as they end. A \omterm{def:b2:threads-and-synchronisation:thread}{thread} per connection is lighter, since a \omterm{def:b2:threads-and-synchronisation:thread}{thread} shares the
\omterm{def:b2:processes-and-system-calls:process}{process}'s memory and is cheaper to start:

\omcode{code/bachelor-2/16-sockets/listings/servers.c}{53}{64}{A \omterm{def:b2:threads-and-synchronisation:thread}{thread} per connection, detached}

\begin{definition}[Event loop]\label{def:b2:sockets:event-loop}
An \emph{event loop}\index{event loop} serves many connections from one \omterm{def:b2:threads-and-synchronisation:thread}{thread}: it asks the system, with
\code{poll(fds, n, timeout)}, which of a set of descriptors are ready (a \code{recv} or an \code{accept} on them would
not wait), handles each ready one briefly, and asks again. A connection's state between two events, such as a request
received in part, is kept in a structure of the program.
\end{definition}

\omcode{code/bachelor-2/16-sockets/listings/servers.c}{69}{104}{An \omterm{def:b2:sockets:event-loop}{event loop}: one \omterm{def:b2:threads-and-synchronisation:thread}{thread}, one \code{poll}, every
connection's state in \code{st}}

\begin{omfigure}
\begin{tikzpicture}[font=\scriptsize, d/.style={draw, minimum width=1.3cm, minimum height=5mm, inner sep=1pt},
  w/.style={-{Stealth[length=1.8mm]}}]
\node[draw, rounded corners=3pt, fill=omThm!12, minimum width=2.2cm, minimum height=9mm, align=center] (p) at (0,0)
{\code{poll}\\ waits};
\node[d, fill=omDef!25] (l) at (4,1.2) {listening};
\node[d, fill=omDef!25] (a) at (4,0.4) {client A};
\node[d] (b) at (4,-0.4) {client B};
\node[d, fill=omDef!25] (c) at (4,-1.2) {client C};
\foreach \x in {l,a,b,c} \draw[w, gray] (p.east) -- (\x.west);
\node[anchor=west, align=left] at (5.0,1.2) {ready: \code{accept}, add the new socket};
\node[anchor=west, align=left] at (5.0,0.4) {ready: \code{recv}; a whole request? answer, close};
\node[anchor=west, align=left] at (5.0,-0.4) {not ready: skipped};
\node[anchor=west, align=left] at (5.0,-1.2) {ready: \code{recv}; part of a request: keep it};
\draw[w] (4,-1.6) -- ++(0,-0.35) -| node[pos=0.25, below] {then \code{poll} again} (p.south);
\end{tikzpicture}
\par\smallskip{\small An \omterm{def:b2:sockets:event-loop}{event loop}. One call reports every ready descriptor; each is served without waiting, and
the loop returns to \code{poll}. Nothing blocks, so one \omterm{def:b2:threads-and-synchronisation:thread}{thread} keeps every connection moving.}
\end{omfigure}

Which is fastest? \code{loadtest.c} starts each server in a \omterm{def:b2:processes-and-system-calls:tree}{child process} and sends it 4000 connections from 8
client \omterm{def:b2:threads-and-synchronisation:thread}{threads}, each connection with one request:

\omoutput{code/bachelor-2/16-sockets/outputs/load.txt}{One run on the author's laptop: 4000 connections per
design.}\label{out:b2:sockets:load}

\begin{omfigure}
\begin{tikzpicture}
\begin{axis}[omrecord, width=0.68\linewidth, height=4.8cm, ybar, bar width=18pt, ymin=0,
    symbolic x coords={process,thread,poll}, xtick=data, enlarge x limits=0.3,
    xticklabels={process per connection, thread per connection, event loop},
    xticklabel style={font=\scriptsize}, ylabel={requests per second},
    nodes near coords, nodes near coords style={font=\scriptsize}, scaled y ticks=false,
    yticklabel style={/pgf/number format/1000 sep={\,}}]
  \addplot[fill=omDef!40, draw=omDef] table[x=design, y=rate]
    {code/bachelor-2/16-sockets/figures/designs.dat};
\end{axis}
\end{tikzpicture}
\par\smallskip{\small The rates of \cref{out:b2:sockets:load}. Each connection costs a \code{fork}, a \omterm{def:b2:threads-and-synchronisation:thread}{thread}, or
nothing extra; with one short request per connection, that cost is most of the work.}
\end{omfigure}

The order follows what each design pays per connection: a \omterm{def:b2:processes-and-system-calls:process}{process} (the copy of the page tables, the exit, the
reaping), a \omterm{def:b2:threads-and-synchronisation:thread}{thread} (a stack and a scheduling entry), or an entry in an array. With long connections that carry many
requests, the gap narrows; with thousands of idle connections, an \omterm{def:b2:sockets:event-loop}{event loop} needs no \omterm{def:b2:threads-and-synchronisation:thread}{thread} for each, which is why
busy servers are built on it (\cref{exo:b2:sockets:11}). Its weakness is that one slow handler stops every client: a
handler must never wait.

\begin{inpractice}[Ports and restarts]
A server restarted at once often fails with ``Address already in use'': its old connections linger for a while in the
\code{TIME_WAIT} state, to absorb late packets. The option \code{SO_REUSEADDR}, set before \code{bind}, allows the
bind anyway, and every server of this chapter sets it (\cref{exo:b2:sockets:12}). Ports below 1024 are reserved to
privileged programs; servers under test use high ports. % ledger: b2:ephemeral-range
\begin{dated}{2026}
On Linux, a client's own port is chosen from 32\,768 to 60\,999 by default; the load test's 12\,000 short
connections borrow a good share of them for a minute.
\end{dated}
\end{inpractice}

\section{Exercises}\label{sec:b2:sockets:exercises}

\begin{exercise}[$\star$]\label{exo:b2:sockets:1}
A server sends ``ab'', then ``cd'', then closes. The client reads with a buffer of 3 bytes and prints each piece in
brackets. List every output it may print.
\omcode{code/bachelor-2/16-sockets/listings/exo_abcd.c}{9}{23}{Both sides}
\end{exercise}

\begin{exercise}[$\star$]\label{exo:b2:sockets:2}
In \code{wrong_port.c}, compute the port the server really listens on, from the bytes of 44164, and explain why
\code{bind} succeeded.
\end{exercise}

\begin{exercise}[$\star$]\label{exo:b2:sockets:3}
Write \code{recv_all(fd, buf, n)} as specified by \cref{prop:b2:sockets:recv-all}, and \code{send_all}.
\end{exercise}

\begin{exercise}[$\star$]\label{exo:b2:sockets:4}
Write in hexadecimal the bytes on the wire for the framed request \code{street Mill Lane}. How many bytes does a
\code{recv_msg} read, in how many calls to \code{recv_all}?
\end{exercise}

\begin{exercise}[$\star\star$]\label{exo:b2:sockets:5}
What does this program print, and why? What happens to the other 90 bytes?
\omcode{code/bachelor-2/16-sockets/listings/exo_truncate.c}{19}{26}{A datagram of 100 bytes, a buffer of 10}
\end{exercise}

\begin{exercise}[$\star\star$]\label{exo:b2:sockets:6}
This echo server spins after its client leaves. Explain the output and fix the loop.
\omcode{code/bachelor-2/16-sockets/listings/exo_loop.c}{25}{30}{The server's loop}
\omoutput{code/bachelor-2/16-sockets/outputs/exo_loop.txt}{The loop, stopped by its guard.}
\end{exercise}

\begin{exercise}[$\star\star$]\label{exo:b2:sockets:7}
Write a line reader, \code{lr_line(r, out, cap)}, that returns the next line of a stream without its newline, using a
buffer of its own so that it calls \code{recv} with large sizes and never reads past what it needs to keep.
\end{exercise}

\begin{exercise}[$\star\star$]\label{exo:b2:sockets:8}
Prove \cref{prop:b2:sockets:prefix-free} by induction on $k$. Show by an example that with a delimiter, and a message
that contains the delimiter, decoding is not unique.
\end{exercise}

\begin{exercise}[$\star\star$]\label{exo:b2:sockets:9}
Write a server with a \omterm{def:b2:threads-and-synchronisation:thread}{thread} per client, in which each line a client sends is answered with the number of lines
received from all clients so far.
\end{exercise}

\begin{exercise}[$\star\star\star$]\label{exo:b2:sockets:10}
Write a chat server with \code{poll}: each line a client sends is sent to every other connected client.
\end{exercise}

\begin{exercise}[$\star\star\star$]\label{exo:b2:sockets:11}
Model the time per connection of each design as a fixed cost $f$ plus the cost $r$ of a request, the same for all.
From \cref{out:b2:sockets:load}, estimate $f$ for a \omterm{def:b2:processes-and-system-calls:process}{process} and for a \omterm{def:b2:threads-and-synchronisation:thread}{thread}, taking $f = 0$ for the \omterm{def:b2:sockets:event-loop}{event loop}. How
many requests per connection would make the \omterm{def:b2:processes-and-system-calls:process}{process} design reach 90~\% of the \omterm{def:b2:sockets:event-loop}{event loop}'s rate?
\end{exercise}

\begin{exercise}[$\star\star\star$]\label{exo:b2:sockets:12}
Explain the output of this experiment, in which a server closes a connection first, ends, and then a new server binds
the same port without and with \code{SO_REUSEADDR}.
\omoutput{code/bachelor-2/16-sockets/outputs/reuse.txt}{Two binds, right after a server
closed.}
\end{exercise}

\section{Problem: The town directory service}\label{sec:b2:sockets:problem}

\begin{problem}[{Weekend problem --- The town directory service}]\label{pb:b2:sockets:1}
The problem follows the directory service of \cref{out:b2:sockets:directory} from its protocol to its load test.

\needspace{4\baselineskip}\medskip\noindent\textbf{Part I --- The protocol.}\par\nopagebreak
\begin{enumerate}
\item Write the bytes of the framed request \code{who 1}, as the server receives them.
\item What is the largest message a 4-byte length can describe? Why does \code{recv_msg} take a capacity?
\item A client sends \code{who 1} and \code{who 2} on one connection, back to back. Can the server's first
\code{recv} return parts of both? Why does framing make it harmless?
\item The answer to \code{who 1} is 39 bytes long. Write its first four bytes.
\item The Python client writes \code{struct.pack("!I", n)}. What would \code{"<I"} send instead, and what would the C
server then read as the length of \code{who 1}?
\item Why is the reply for an unknown request a message, rather than a closed connection?
\end{enumerate}

\needspace{4\baselineskip}\medskip\noindent\textbf{Part II --- The server.}\par\nopagebreak
\begin{enumerate}[resume, beginpenalty=10000]
\item List the \omterm{def:b2:processes-and-system-calls:syscall}{system calls} of \code{tcp_listen}, in order, and say what each one does.
\item What does \code{accept} return, and why is it a new descriptor rather than the listening one?
\item In \code{handle}, what happens if a client connects and sends nothing, with each design?
\item Use \cref{prop:b2:sockets:recv-all} to show that \code{recv_msg} returns a whole message or $-1$.
\item \code{send_all} passes \code{MSG_NOSIGNAL}. What would happen to the server without it, when a client closes
before its answer is sent (\cref{ch:b2:files-pipes-signals})?
\item Why does the demo server serve exactly four connections and stop?
\end{enumerate}

\needspace{4\baselineskip}\medskip\noindent\textbf{Part III --- Concurrency, three ways.}\par\nopagebreak
\begin{enumerate}[resume, beginpenalty=10000]
\item In \code{serve_fork}, why must the parent close \code{c}? What would happen to the clients otherwise?
\item In \code{serve_threads}, why is the descriptor passed by value, cast to a pointer, and not as \code{&c}?
\item The \omterm{def:b2:threads-and-synchronisation:thread}{threads} share \code{directory_answer}'s tables. Why is no \omterm{def:b2:threads-and-synchronisation:lock}{lock} needed?
\item In \code{serve_poll}, what happens when a request arrives in two pieces?
\item Why does \code{serve_poll} walk its array from the end when it removes entries?
\item Which design survives a handler that crashes with a segmentation fault, and which do not?
\end{enumerate}

\needspace{4\baselineskip}\medskip\noindent\textbf{Part IV --- The load test.}\par\nopagebreak
\begin{enumerate}[resume, beginpenalty=10000]
\item How many connections does each design receive, and from how many \omterm{def:b2:threads-and-synchronisation:thread}{threads}?
\item Why does the load test run each server in a \omterm{def:b2:processes-and-system-calls:tree}{child process} rather than a \omterm{def:b2:threads-and-synchronisation:thread}{thread} of its own?
\item Why is the measured time taken on the clients' side?
\item Order the designs by \cref{out:b2:sockets:load} and give the ratio of the fastest to the slowest.
\item Why might the order change if each connection carried a thousand requests?
\item State the result: requests per second on loopback for the three server designs, on the laptop.
\end{enumerate}
\end{problem}
